\chapter{Scoring Engine and Evaluation}
\label{chapter-7-scoring-and-classification}

\section{The Orchestration of
Intelligence}\label{the-orchestration-of-intelligence}

A primary challenge in engineering a multi-modal threat detection
platform is the synthesis of disparate analytical signals into a
unified, actionable verdict. PhishGuard achieves this through a
sophisticated scoring and classification engine that seamlessly bridges
the continuous probabilistic output of the XGBoost machine learning
model with discrete, rule-based heuristic indicators.

This architecture ensures that the system is not only highly accurate
but also highly transparent, providing Security Operations Center (SOC)
analysts and end-users with explicit, human-readable explanations for
every classification.

\section{The ML Predictor
Architecture}\label{the-ml-predictor-architecture}

The core inference engine is housed within the
\texttt{backend/ml/predictor.py} module. To guarantee high-throughput
and sub-millisecond inference times during production deployment, the
\texttt{PhishingPredictor} class is implemented utilizing a thread-safe
Singleton design pattern.

\subsection{Thread-Safe
Initialization}\label{thread-safe-initialization}

During the FastAPI application startup sequence, the
\texttt{PhishingPredictor} invokes its \texttt{\_\_new\_\_} method,
locking the execution thread (\texttt{threading.Lock()}) to instantiate
the XGBoost model exactly once. The model weights are loaded into memory
from \texttt{models/phishingModel.json}, alongside the strict
21-dimensional feature schema defined in \texttt{modelMetadata.json}.
Once initialized, the model state becomes immutable, allowing the
asynchronous web server to process concurrent inference requests across
multiple threads without encountering race conditions or requiring
continuous disk I/O operations.

\subsection{Inference Execution}\label{inference-execution}

When a URL analysis request is routed to the predictor, the
21-dimensional \texttt{FeatureSet} is cast into a strongly typed
\texttt{numpy.ndarray} (\texttt{dtype=np.float64}). The
\texttt{predict()} method invokes the XGBoost \texttt{predict\_proba}
function, returning a continuous probability bounded between 0.0 and
1.0, representing the mathematical likelihood that the input belongs to
the malicious class.

\section{The Phishing Scorer
Module}\label{the-phishing-scorer-module}

While the XGBoost probability serves as the primary ground truth for URL
inputs, it operates as an opaque numerical value. To resolve the
black-box interpretability crisis, the \texttt{backend/ml/scorer.py}
module introduces a parallel heuristic scoring engine that
reverse-engineers the threat context.

\subsection{Explanatory Heuristics}\label{explanatory-heuristics}

Regardless of whether the XGBoost model is active, the
\texttt{PhishingScorer} executes three discrete heuristic calculators to
generate the human-readable \texttt{factors} array that will be
presented to the user: 1. \textbf{\texttt{calculateUrlStructureScore}:}
Applies deterministic penalty weights to structural anomalies. For
example, a URL length exceeding 75 characters generates a scaling
penalty (\texttt{min((urlLength\ -\ 75)\ /\ 100,\ 0.3)}), while the
presence of raw IP addresses or Suspicious TLDs adds fixed penalties. 2.
\textbf{\texttt{calculateOsintScore}:} Evaluates the infrastructure
context. If \texttt{isNewlyRegistered} is true (domain age $<${}
30 days), a 0.35 penalty is recorded. Detection in third-party
blacklists adds up to a 0.5 penalty. 3.
\textbf{\texttt{calculateFeatureScore}:} Analyzes holistic combinations
that indicate sophisticated evasion, such as a newly registered domain
combined with the presence of suspicious keywords in the URI string
(\texttt{score\ +=\ 0.3}).

\subsection{Graceful Heuristic
Fallback}\label{graceful-heuristic-fallback}

In the event of a catastrophic failure where the compiled XGBoost model
cannot be loaded into memory, the \texttt{PhishingScorer} implements a
graceful degradation protocol. Instead of halting the application, it
relies entirely on the heuristic calculators, combining their raw
outputs using the strict constants defined in the
\texttt{ScoringWeights} dataclass: - \textbf{URL Structure:} 25\% weight
- \textbf{OSINT Derived:} 35\% weight - \textbf{Feature Based:} 40\%
weight

This fallback mechanism ensures the platform remains operational and
capable of identifying overt threats even during degraded system states.

\section{Risk Level Determination}\label{risk-level-determination}

The final continuous score (whether derived from the XGBoost probability
or the heuristic fallback) is mapped to discrete, actionable categories
using the strict boundaries defined in the \texttt{RISK\_THRESHOLDS}
dictionary:

\begin{itemize}
\item
  \textbf{Safe (Score $<${} 0.30):} The domain exhibits normal
  structural characteristics, extensive DNS history, and clean
  reputation metrics. No action required.
\item
  \textbf{Suspicious (Score $<${} 0.50):} The domain exhibits some
  characteristics commonly associated with phishing infrastructure,
  requiring user caution and verification. Minor anomalies may include
  newly registered domains or non-standard configurations.
\item
  \textbf{Dangerous (Score $<${} 0.70):} Strong indications of
  malicious intent, such as obscured IPs, known bad TLDs, missing MX
  records on a young domain, or multiple risk indicators. The system
  recommends not interacting with the content.
\item
  \textbf{Critical (Score $\ge$ 0.70):} Definitive malicious
  classification. The XGBoost model expresses high confidence in the
  threat, or the domain has been explicitly flagged by external threat
  intelligence APIs. Immediate blocking and reporting recommended.
\end{itemize}

\subsection{Reason Prioritization}\label{reason-prioritization}

To prevent alert fatigue and ensure the user interface remains
uncluttered, the \texttt{\_prioritizeReasons} method deduplicates and
sorts the aggregated heuristic factors. It employs a keyword-based
sorting algorithm, ensuring that highly critical warnings (e.g.,
``malicious,'' ``ip address,'' ``newly registered'') are prioritized at
the top of the array, while lower-severity structural warnings are
pushed down or truncated. The final output is strictly limited to the
top 10 most relevant reasons.



\section{Evaluation Methodology}\label{evaluation-methodology}

The efficacy of the PhishGuard machine learning architecture was
rigorously evaluated using a holdout methodology to prevent data leakage
and ensure generalized performance on unseen data. The dataset,
comprising diverse phishing and legitimate Uniform Resource Locators
(URLs), was systematically split into training and testing partitions.

The primary training corpus consisted of 23,374 samples
(\texttt{trainSamples:\ 23,374}). To accurately assess the model's
predictive capability in a simulated real-world environment, a strictly
held-out test set comprising 5,009 samples
(\texttt{testSamples:\ 5,009}) was established. This test set maintained
an near-perfect class balance, containing 2,505 legitimate instances and
2,504 phishing instances, ensuring that evaluation metrics were not
artificially skewed by class imbalance.

\section{Hyperparameter
Optimization}\label{hyperparameter-optimization}

Prior to final evaluation, the XGBoost classifier underwent extensive
hyperparameter tuning utilizing the Optuna optimization framework. The
optimization process executed 50 distinct trials (\texttt{nTrials:\ 50})
employing 5-fold cross-validation (\texttt{nCvFolds:\ 5}) to iteratively
search the hyperparameter space for the configuration that maximized the
Area Under the Receiver Operating Characteristic Curve (ROC AUC).

The optimal configuration was identified at Trial 43
(\texttt{bestTrialNumber:\ 43}), yielding a cross-validated AUC of
0.9943. The resulting best hyperparameters significantly constrained the
model's complexity to prevent overfitting while maintaining high
predictive capacity: - \textbf{Maximum Depth (\texttt{max\_depth}):} 7 -
\textbf{Number of Estimators (\texttt{n\_estimators}):} 700 -
\textbf{Learning Rate (\texttt{learning\_rate}):} \textasciitilde0.177 -
\textbf{Subsample Ratio (\texttt{subsample}):} \textasciitilde0.945 -
\textbf{Column Subsample by Tree (\texttt{colsample\_bytree}):}
\textasciitilde0.873 - \textbf{Gamma (\texttt{gamma}):}
\textasciitilde0.198

\section{Empirical Performance
Metrics}\label{empirical-performance-metrics}

The fully optimized XGBoost model, utilizing the complete 21-dimensional
feature vector (17 URL structural features and 4 OSINT-derived
features), demonstrated exceptional performance on the held-out test set
of 5,009 samples. 

As summarized in Table~\ref{tab:performance_metrics}, the model achieved an aggregate \textbf{Accuracy of 96.45\%}
(\texttt{0.96446}), indicating robust overall correctness. However, in
the context of threat detection, precision and recall offer more nuanced
insights into operational viability.

\begin{table}[htbp]
\centering
\caption{PhishGuard XGBoost Model Performance Metrics (Test Set: N=5,009)}
\label{tab:performance_metrics}
\renewcommand{\arraystretch}{1.3}
\begin{tabular}{@{}ll@{}}
\toprule
\textbf{Metric} & \textbf{Value (\%)} \\
\midrule
Accuracy & 96.45\% \\
Precision & 97.86\% \\
Recall (Sensitivity) & 94.97\% \\
F1-Score & 96.39\% \\
ROC AUC & 99.41\% \\
PR AUC & 99.48\% \\
\bottomrule
\end{tabular}
\end{table}

The system achieved a \textbf{Precision of 97.86\%} (\texttt{0.97860}),
demonstrating a low false-positive rate---a critical metric for
minimizing alert fatigue among end-users. Correspondingly, the
\textbf{Recall reached 94.97\%} (\texttt{0.94968}), signifying the
model's strong capability to successfully identify true phishing threats
within the dataset. The harmonic mean of these two metrics resulted in
an \textbf{F1-Score of 96.39\%} (\texttt{0.96392}).

Furthermore, the model exhibited outstanding discriminative ability
across various classification thresholds, achieving a \textbf{ROC AUC of
99.41\%} (\texttt{0.99408}) and a Precision-Recall Area Under Curve
(\textbf{PR AUC}) of \textbf{99.48\%} (\texttt{0.99479}).

\subsection{Error Analysis and Confusion
Matrix}\label{error-analysis-and-confusion-matrix}

A granular analysis of the classification errors on the test set reveals
the model's operational tendencies. Table~\ref{tab:confusion_matrix} presents the exact distribution of the 5,009 test predictions.

\begin{table}[htbp]
\centering
\caption{Confusion Matrix for the Holdout Test Set}
\label{tab:confusion_matrix}
\renewcommand{\arraystretch}{1.5}
\begin{tabular}{cc|cc}
\multicolumn{2}{c}{} & \multicolumn{2}{c}{\textbf{Predicted Class}} \\
\multicolumn{2}{c|}{} & \textbf{Legitimate (0)} & \textbf{Phishing (1)} \\
\hline
\multirow{2}{*}{\textbf{Actual}} & \textbf{Legitimate (0)} & 2,453 (TN) & 52 (FP) \\
& \textbf{Phishing (1)} & 126 (FN) & 2,378 (TP) \\
\end{tabular}
\end{table}

The disparity between False Positives (52) and False Negatives (126)
underscores the model's high-precision orientation. While the system is
highly reliable when it issues a warning, a small subset of
sophisticated phishing URLs successfully evaded the structural and OSINT
feature checks.

\section{Feature Explainability and SHAP
Analysis}\label{feature-explainability-and-shap-analysis}

To demystify the XGBoost model's decision-making process and move beyond
``black-box'' predictions, two complementary explainability methods were
employed: XGBoost's native \emph{gain-based} feature importance and
SHapley Additive exPlanations (SHAP) via the \texttt{TreeExplainer}.
These methods measure feature importance differently---gain importance
quantifies the total split-improvement across all trees, while SHAP
assigns each feature a marginal contribution grounded in cooperative
game theory---and their rankings need not coincide.

\paragraph{XGBoost Gain-Based Importance}

The native gain importance (\texttt{evaluation\_report.json}), which
measures the average improvement in the loss function attributable to
each feature across all tree splits, revealed the following hierarchy.
As summarized in Table~\ref{tab:gain_importance}, the presence of HTTPS
(\texttt{isHttps}) emerged as the most influential single feature
(accounting for roughly 33.4\% of the total gain), followed by the
validity of the domain's DNS configuration (\texttt{hasValidDns},
12.6\%). Structural anomalies such as \texttt{specialCharCount} (8.5\%)
and \texttt{pathDepth} (7.4\%) also demonstrated significant predictive
power. Notably, three of the four OSINT features (\texttt{usesCdn},
\texttt{hasValidMx}, \texttt{hasValidDns}) ranked in the lower half of
the gain-based hierarchy, and four features
(\texttt{hasUnderscoreInDomain}, \texttt{hasDoubleSlash},
\texttt{hasIpAddress}, \texttt{hasPortNumber}) registered zero gain,
indicating they were never selected for a beneficial split.

\begin{table}[htbp]
\centering
\caption{Top-10 Features by XGBoost Gain Importance}
\label{tab:gain_importance}
\renewcommand{\arraystretch}{1.2}
\begin{tabular}{@{}rllr@{}}
\toprule
\textbf{Rank} & \textbf{Feature} & \textbf{Category} & \textbf{Gain (\%)} \\
\midrule
1 & \texttt{isHttps} & URL & 33.4 \\
2 & \texttt{hasValidDns} & OSINT & 12.6 \\
3 & \texttt{specialCharCount} & URL & 8.5 \\
4 & \texttt{pathDepth} & URL & 7.4 \\
5 & \texttt{hasAtSymbol} & URL & 5.4 \\
6 & \texttt{hasEncodedChars} & URL & 5.2 \\
7 & \texttt{subdomainCount} & URL & 4.5 \\
8 & \texttt{queryParamCount} & URL & 4.3 \\
9 & \texttt{domainLength} & URL & 4.3 \\
10 & \texttt{dnsRecordCount} & OSINT & 4.0 \\
\bottomrule
\end{tabular}
\end{table}

\paragraph{SHAP-Based Importance}

The SHAP \texttt{TreeExplainer} computes the exact marginal contribution
of each feature to individual predictions using Shapley values from
cooperative game theory \cite{lundberg2017}. Unlike gain importance,
which is aggregated at the tree level, SHAP values provide per-instance
explanations that are locally consistent. The mean absolute SHAP values
(\texttt{ablation\_report.json}) revealed a \emph{different} ranking:
\texttt{urlLength} dominated with a mean $|\Phi|$ of 3.12, followed by
\texttt{pathDepth} (1.63) and \texttt{dnsRecordCount} (0.91). The
\texttt{isHttps} feature, which ranked first in gain importance, dropped
to mid-rank under SHAP (0.38), reflecting the fact that while HTTPS
partitions the data effectively in tree splits, its per-prediction
marginal contribution is comparatively modest once other features are
accounted for.

This divergence between the two methods is not contradictory but
instructive: gain importance highlights \emph{which features the trees
rely on structurally}, while SHAP reveals \emph{which features
materially shift individual predictions}. Both perspectives are valuable
for model auditing---gain for understanding the learned tree structure,
and SHAP for delivering trustworthy, instance-level explanations to
end-users.

\subsection{OSINT Ablation Study}\label{osint-ablation-study}

A critical objective of the PhishGuard architecture was to evaluate the
supplemental value of Open-Source Intelligence (OSINT) data when
combined with traditional URL lexical analysis. To quantify this, an
ablation study was conducted: the XGBoost classifier was trained and
evaluated under two configurations---(a) the full 21-feature vector (17
URL + 4 OSINT) and (b) a reduced 17-feature vector (URL-only)---with all
other parameters held constant.

\paragraph{SHAP Contribution}

The global SHAP contribution analysis (\texttt{ablation\_report.json})
calculated the mean absolute SHAP values for the two distinct feature
subsets. The analysis determined that the 17 URL structural features
account for \textbf{85.64\%} of the model's total SHAP-based feature
importance. The 4 dynamic OSINT features
(\texttt{usesCdn}, \texttt{dnsRecordCount},
\texttt{hasValidDns}, \texttt{hasValidMx}) contributed the remaining
\textbf{14.36\%}.

\paragraph{Predictive Performance Comparison}

Table~\ref{tab:ablation} compares the two configurations across standard
classification metrics. Notably, the URL-only model achieved marginally
\emph{higher} raw accuracy than the OSINT-augmented model
($98.93\%$ vs.\ $98.40\%$, a difference of $-0.53$ percentage points
when OSINT is added). This difference is within the bounds of
cross-validation variance (CV AUC standard deviations of 0.0012 and
0.0006 respectively) and should not be interpreted as OSINT
\emph{hurting} the model; rather, it indicates that the 17 lexical
features were already sufficiently discriminative on this particular
balanced dataset to classify the targets independently.

\begin{table}[htbp]
\centering
\caption{OSINT Ablation Study: URL-Only vs.\ URL+OSINT}
\label{tab:ablation}
\renewcommand{\arraystretch}{1.3}
\begin{tabular}{@{}lrrr@{}}
\toprule
\textbf{Metric} & \textbf{URL-Only (17 feat.)} & \textbf{URL+OSINT (21 feat.)} & \textbf{Difference} \\
\midrule
Accuracy & 98.93\% & 98.40\% & $-0.53$ pp \\
Precision & 99.78\% & 99.13\% & $-0.65$ pp \\
Recall & 98.08\% & 97.65\% & $-0.43$ pp \\
F1-Score & 98.92\% & 98.39\% & $-0.54$ pp \\
ROC AUC & 99.80\% & 99.77\% & $-0.04$ pp \\
\bottomrule
\end{tabular}
\end{table}

\paragraph{Interpretation}

The inclusion of OSINT features, while contributing 14.36\% to the
model's SHAP explanations and providing critical human-readable context
for the heuristic Scorer module, did not yield a raw accuracy
improvement on this dataset. This finding highlights a crucial
distinction in applied cybersecurity ML: the value of OSINT integration
is \emph{not} primarily in boosting the classifier's raw predictive
metric but in enabling \emph{explainable} threat intelligence. The OSINT
features supply infrastructure-level context (domain age, DNS validity,
CDN usage, MX record presence) that the \texttt{PhishingScorer} module
translates into actionable, human-readable justifications---context that
purely lexical models cannot provide. A security analyst informed that a
domain was ``registered 3 days ago and lacks valid MX records'' gains
far more actionable intelligence than a bare 96\% phishing probability
score, even if the latter is marginally more accurate. This trade-off
between raw predictive performance and explainable, context-rich
intelligence is a deliberate architectural decision discussed further in
the conclusion (Chapter~\ref{chapter-11-discussion}).

